\documentclass[11pt,a4paper]{article}

\usepackage[T1]{fontenc}
\usepackage[utf8]{inputenc} % pdfLaTeX
\usepackage{lmodern}
\usepackage{microtype}
\usepackage{amsmath,amssymb}
\usepackage[a4paper,margin=2.2cm]{geometry}
\usepackage{enumitem}
\usepackage[hidelinks]{hyperref}

% Headings already include manual numbering (e.g. "1. ..."), so disable LaTeX numbering:
\setcounter{secnumdepth}{-1}

\title{Oxford Handbook of Anaesthesia (3e, 2011)\\Chapter 15: Thoracic surgery\\Structured study notes}
\author{}
\date{}

\begin{document}
\maketitle

\noindent\textbf{Source:} Keith Allman \& Iain Wilson (eds.) --- \emph{Oxford Handbook of Anaesthesia} (3rd ed., 2011), Chapter 15 ``Thoracic surgery''.

\tableofcontents

\section{1. General principles}

\subsection{1.1 What thoracic anaesthesia demands}
\begin{itemize}[leftmargin=*, itemsep=2pt]
  \item Ability to control ventilation of \textbf{each lung independently}.
  \item Skilled management of the \textbf{shared airway/lung} with the surgeon.
  \item Clear understanding of the planned surgery.
  \item \textbf{Communication} between surgeon and anaesthetist is essential.
\end{itemize}

\subsection{1.2 Typical patient profile and implications}
Thoracic surgical patients are commonly \textbf{older} and \textbf{less fit} than other surgical populations. Associated problems often include long-term smoking, bronchial carcinoma, pleural effusion, empyema, oesophageal obstruction, cachexia \(\rightarrow\) reduced \textbf{cardiorespiratory reserve}.

\subsection{1.3 General considerations}
\begin{itemize}[leftmargin=*, itemsep=2pt]
  \item Discuss the procedure and potential problems with the surgeon.
  \item Optimise lung function before elective surgery:
  \begin{itemize}[leftmargin=*, itemsep=2pt]
    \item Smoking cessation.
    \item Pre-op physiotherapy + incentive spirometry.
    \item Optimise bronchodilators; consider a course of oral steroids.
  \end{itemize}
  \item Position: most procedures use \textbf{lateral decubitus}, table ``broken'' to separate ribs.
  \item Avoid postoperative mechanical ventilation if possible:
  \begin{itemize}[leftmargin=*, itemsep=2pt]
    \item Stresses pulmonary suture lines.
    \item Increases air leaks + risk of chest infection.
  \end{itemize}
  \item Minimise postoperative respiratory dysfunction: good analgesia + physiotherapy.
  \item Postoperative oxygen therapy commonly prescribed to compensate for increased V/Q mismatch:
  \begin{itemize}[leftmargin=*, itemsep=2pt]
    \item After pulmonary surgery: warmed humidified \(40\%\) \(\mathrm{O_2}\) by facemask recommended.
    \item Nasal cannulae (e.g. 3 L/min) better tolerated and satisfactory for most other patients.
  \end{itemize}
\end{itemize}

\subsection{1.4 Preoperative assessment (emphasis on lung resection)}
\begin{itemize}[leftmargin=*, itemsep=2pt]
  \item Standard assessment with particular emphasis on \textbf{cardiorespiratory reserve}.
  \item Review latest CXR and CT for:
  \begin{itemize}[leftmargin=*, itemsep=2pt]
    \item Airway obstruction.
    \item Tracheal/carinal distortion or compression (impacts DLT placement).
    \item Tumour relations: fissures/chest wall/major vessels (surgical implications).
  \end{itemize}
  \item Patients with significant cardiac disease are high risk.
\end{itemize}

\subsubsection{1.4.1 Lung resection: practical risk stratification}
\begin{itemize}[leftmargin=*, itemsep=2pt]
  \item Clinically fit, good exercise tolerance, normal spirometry \(\rightarrow\) generally acceptable.
  \item Major medical problems + minimal exercise capacity + grossly impaired PFTs \(\rightarrow\) high risk; consider alternatives.
  \item Reduced exercise capacity (e.g. dyspnoea climbing 2 flights of stairs) + abnormal spirometry \(\pm\) co-morbidity \(\rightarrow\) careful evaluation.
\end{itemize}

\subsubsection{1.4.2 PFTs: interpretation}
\begin{itemize}[leftmargin=*, itemsep=2pt]
  \item Spirometry reflects ``bellows'' function.
  \item Diffusion capacity (e.g. DLCO/transfer factor) reflects ability to transfer \(\mathrm{O_2}\) to the circulation.
  \item Diffuse alveolar disease may show \textbf{severely impaired gas transfer} with relatively normal spirometry.
\end{itemize}

\subsubsection{1.4.3 Commonly accepted minimum FEV\(_1\) thresholds (\% predicted)}
\begin{itemize}[leftmargin=*, itemsep=2pt]
  \item Pneumonectomy: \(\textgreater 55\%\)
  \item Lobectomy: \(\textgreater 40\%\)
  \item Wedge resection: \(\textgreater 35\%\)
\end{itemize}

\section{2. Analgesia}

\subsection{2.1 Why it matters}
\begin{itemize}[leftmargin=*, itemsep=2pt]
  \item Thoracotomy incisions are \textbf{extremely painful}.
  \item Poor pain control increases stress response and impairs mobilisation/respiration \(\rightarrow\) more respiratory complications.
  \item Chronic post-thoracic surgery pain syndrome occurs in \textbf{25--60\%}.
\end{itemize}

\subsection{2.2 Recommended analgesic strategy (multimodal)}
Combine:
\begin{itemize}[leftmargin=*, itemsep=2pt]
  \item Paracetamol.
  \item NSAIDs.
  \item Regional block.
  \item Intraoperative opioids.
  \item Regular or PCA postoperative analgesia.
\end{itemize}

\subsection{2.3 Procedure-linked options}
\begin{itemize}[leftmargin=*, itemsep=2pt]
  \item Thoracoscopic procedures:
  \begin{itemize}[leftmargin=*, itemsep=2pt]
    \item Intercostal nerve blocks or percutaneous paravertebral block.
    \item Oral or PCA opioid post-op.
  \end{itemize}
  \item Thoracotomy / thoracoabdominal incisions:
  \begin{itemize}[leftmargin=*, itemsep=2pt]
    \item Unless contraindicated: \textbf{continuous thoracic epidural} OR \textbf{continuous paravertebral} analgesia.
    \item Epidural and paravertebral have \textbf{equivalent} analgesic efficacy and similar effects on stress response/respiratory function.
    \item Paravertebral tends to have fewer adverse events: hypotension, urinary retention, PONV.
  \end{itemize}
\end{itemize}

\subsection{2.4 Practical points}
\begin{itemize}[leftmargin=*, itemsep=2pt]
  \item Match block level to incision: usually \textbf{T5/6 or T6/7}.
  \item Thoracic epidural: establish cautiously with small incremental boluses (e.g. 3--4 mL of 0.25\% bupivacaine) to avoid severe hypotension from extensive sympathetic block.
  \item Paravertebral:
  \begin{itemize}[leftmargin=*, itemsep=2pt]
    \item Example percutaneous injection: 0.5\% bupivacaine (0.3 mL/kg).
    \item Continuous post-op infusion via surgically placed catheter may be used (concentration stepped down after first 24 h).
  \end{itemize}
\end{itemize}

\section{3. Isolation of the lungs}

\subsection{3.1 Key messages}
\begin{itemize}[leftmargin=*, itemsep=2pt]
  \item Independent ventilation is not always straightforward.
  \item OLV has complications and should be used when benefits outweigh risks.
\end{itemize}

\subsection{3.2 Advantages}
\begin{itemize}[leftmargin=*, itemsep=2pt]
  \item Protects dependent lung from blood/secretions.
  \item Allows independent control of ventilation.
  \item Improves surgical access and reduces lung trauma.
\end{itemize}

\subsection{3.3 Disadvantages}
\begin{itemize}[leftmargin=*, itemsep=2pt]
  \item Creates a shunt \(\rightarrow\) usually hypoxia.
  \item Acute lung injury occurs in a small proportion.
  \item Increased technical + physiological challenge.
\end{itemize}

\subsection{3.4 Indications for separation/isolation}
\begin{itemize}[leftmargin=*, itemsep=2pt]
  \item Prevent contamination in infection, massive pulmonary haemorrhage, or bronchoalveolar lavage.
  \item Control ventilation distribution in massive air leaks or severe unilateral disease (e.g. giant bullae, lung cysts).
  \item Surgical access is a relative indication: if OLV is difficult, discuss with surgeon---retraction may suffice.
\end{itemize}

\subsection{3.5 Techniques}
\begin{itemize}[leftmargin=*, itemsep=2pt]
  \item Double lumen tubes (DLTs): commonest and most versatile.
  \item Bronchial blockers (e.g. Univent, Arndt): useful in experienced hands; especially difficult airway/distorted anatomy/tracheostomy.
  \item DLTs and bronchial blockers are clinically equivalent for provision of OLV.
  \item Single lumen endobronchial tubes are rarely used.
\end{itemize}

\subsection{3.6 DLTs: selection and placement essentials}
\begin{itemize}[leftmargin=*, itemsep=2pt]
  \item Disposable PVC DLTs widely used.
  \item Right- vs left-sided DLTs:
  \begin{itemize}[leftmargin=*, itemsep=2pt]
    \item Right-sided tubes have a portal to ventilate the right upper lobe.
    \item Common practice: choose \textbf{left-sided} unless proximal left lobar resection/pneumonectomy or abnormal left bronchial anatomy likely.
  \end{itemize}
  \item Size terminology: Charri\`ere (Ch) = external circumference in mm.
  \item DLT lumens are smaller than standard ETTs; bronchoscopy requires a narrow scope (\(\textless 4\) mm), ideally with self-contained light.
\end{itemize}

\subsubsection{3.6.1 Placement technique highlights}
\begin{itemize}[leftmargin=*, itemsep=2pt]
  \item After tip past cords: partially withdraw stylet and rotate tube to correct orientation.
  \item Use the largest DLT that passes easily through the glottis.
  \item Patient height determines insertion depth; avoid advancing too far.
\end{itemize}

\subsubsection{3.6.2 Confirmation of position}
\begin{itemize}[leftmargin=*, itemsep=2pt]
  \item Confirm lung isolation required for surgery.
  \item Confirm ability to ventilate the non-operative lung: clamp operative limb at Y-connector and open sealing cap to allow deflation.
  \item Observe unilateral chest expansion and correlate with pathology and pre-op exam/radiology.
\end{itemize}

\section{4. Management of one lung ventilation (OLV)}

\subsection{4.1 Principle}
OLV inevitably creates a shunt through the non-ventilated lung. Main goal: \textbf{minimise the effects of shunt}.

\subsection{4.2 Initiating OLV (stepwise)}
\begin{itemize}[leftmargin=*, itemsep=2pt]
  \item Before OLV: increase \(\mathrm{FiO_2}\), reduce tidal volume to a protective range, and note resulting airway pressures.
  \item Start OLV:
  \begin{itemize}[leftmargin=*, itemsep=2pt]
    \item Clamp Y-connection to operative (non-dependent) lung.
    \item Open sealing cap on that lumen to allow gas escape.
    \item Expect airway pressure rise by \(\sim 30\)---\(40\%\) if OLV achieved (may not if operative lung already non-functioning pre-op).
  \end{itemize}
\end{itemize}

\subsection{4.3 If airway pressures are excessive or rise abruptly}
\begin{itemize}[leftmargin=*, itemsep=2pt]
  \item Exclude mechanical causes (kinked connector, clamp misplacement).
  \item Exclude DLT malposition/obstruction (ventilating a lobe not lung, sputum plugs, lumen against airway wall).
  \item Adjust tidal volume/ventilation profile to limit Paw (protective strategy).
  \item Monitor SpO\(_2\) and ETCO\(_2\) closely; increase RR if needed to maintain minute ventilation.
  \item Confirm with surgeon: lung collapse (may be slow in obstructive disease); mediastinum not ``sunk'' into dependent hemithorax.
\end{itemize}

\subsection{4.4 Failure to achieve OLV}
\begin{itemize}[leftmargin=*, itemsep=2pt]
  \item Re-check DLT position and cuff inflation.
  \item Confirm which lumen is clamped and which lung is ventilated.
  \item Bronchoscopic reassessment is often required.
\end{itemize}

\subsection{4.5 Management of hypoxia during OLV (escalation)}
\begin{itemize}[leftmargin=*, itemsep=2pt]
  \item Confirm tube position and suction both lumens.
  \item Increase \(\mathrm{FiO_2}\).
  \item Optimise ventilation of dependent lung.
  \item Apply CPAP/\(\mathrm{O_2}\) to operative lung when feasible (balance against surgical requirements).
  \item Consider recruitment manoeuvres and PEEP strategies (tailored).
  \item If severe hypoxaemia persists \(\rightarrow\) return to two-lung ventilation as required.
\end{itemize}

\section{5. Rigid bronchoscopy and stent insertion}

\subsection{5.1 Preoperative}
\begin{itemize}[leftmargin=*, itemsep=2pt]
  \item Check for airway obstruction (stridor, tracheal tumour on CT, foreign body).
  \item Often suitable as day case.
  \item Warn about postoperative coughing, haemoptysis, and suxamethonium myalgia.
  \item Often combined with mediastinoscopy.
  \item Airway unprotected: if aspiration risk, reduce gastric volume/acidity (e.g. PPI/H\(_2\) blocker regimen described in source).
\end{itemize}

\subsection{5.2 Perioperative}
\begin{itemize}[leftmargin=*, itemsep=2pt]
  \item Full preoxygenation.
  \item Confirm surgeon present before induction.
  \item Induction support as described (e.g. benzodiazepine + short-acting opioid).
  \item Optional small ``taming'' dose of non-depolarising relaxant to reduce suxamethonium pains.
  \item Workflow described: induce in anaesthetic room, transfer with facemask, give suxamethonium just prior.
  \item If potential obstruction: inhalational induction in theatre with sevoflurane in \(\mathrm{O_2}\) until airway secure.
  \item Coordinate ventilation with surgical activity.
  \item Monitor for recovery of muscle tone; suction upper airway and confirm adequate power before scope removal.
\end{itemize}

\subsection{5.3 Postoperative}
\begin{itemize}[leftmargin=*, itemsep=2pt]
  \item Turn patient biopsied side down to reduce spillage to normal lung.
  \item Sit fully upright as soon as awake.
  \item Blood clot can cause severe obstruction \(\rightarrow\) be prepared for immediate airway intervention and repeat bronchoscopy.
\end{itemize}

\subsection{5.4 Special considerations}
\begin{itemize}[leftmargin=*, itemsep=2pt]
  \item Very stimulating \(\rightarrow\) marked hypertensive response.
  \item Need profound relaxation + obtund CV responses \textbf{with rapid return} of laryngeal reflexes and spontaneous respiration.
  \item Topical LA to cords may be used, but will not prevent carinal reflexes and may impair postoperative cough.
\end{itemize}

\section{6. Superior/cervical mediastinoscopy}

\subsection{6.1 Preoperative}
\begin{itemize}[leftmargin=*, itemsep=2pt]
  \item Suitable as day case in appropriate patients.
  \item Check for SVC obstruction and tracheal deviation/compression from mediastinal masses.
  \item Often preceded by rigid bronchoscopy (``Bronch \& Med'').
\end{itemize}

\subsection{6.2 Perioperative}
\begin{itemize}[leftmargin=*, itemsep=2pt]
  \item Tape eyes; check connectors (head obscured by drapes).
  \item Boluses of IV opioid for analgesia.
  \item Insert a large-bore cannula in a lower limb vein after induction (bleeding/SVC issues).
  \item Watch for surgical tracheal compression: monitor tidal volume and airway pressures.
  \item Monitor BP in left arm; pulse oximeter on right hand.
\end{itemize}

\subsection{6.3 Postoperative}
Simple analgesia (e.g. paracetamol + NSAID).

\subsection{6.4 Special considerations}
\begin{itemize}[leftmargin=*, itemsep=2pt]
  \item Risk of massive haemorrhage from great vessels (higher with SVC obstruction) \(\rightarrow\) may require immediate sternotomy.
  \item Mediastinoscope may compress brachiocephalic artery \(\rightarrow\) reduced right arm/carotid flow \(\rightarrow\) cerebral ischaemia risk; right-hand oximeter helps detect perfusion compromise.
  \item Mediastinotomy can cause pneumothorax.
\end{itemize}

\section{7. Lung surgery: wedge resection, lobectomy, pneumonectomy}

\subsection{7.1 Preoperative}
\begin{itemize}[leftmargin=*, itemsep=2pt]
  \item Commonest indication: cancer; others include benign tumours, bronchiectasis, TB.
  \item Assess cardiorespiratory reserve and estimate post-resection lung function.
  \item Assess airway for DLT placement.
  \item Plan postoperative analgesia.
\end{itemize}

\subsection{7.2 Perioperative}
\begin{itemize}[leftmargin=*, itemsep=2pt]
  \item Select appropriate DLT; confirm isolation after intubation.
  \item Prefer left-sided tube unless proximal left lobectomy/pneumonectomy or abnormal left bronchial anatomy.
  \item IV access (non-dependent arm) + arterial line (often better in dependent arm).
  \item CVP unreliable in lateral open chest; central lines not routine (may be needed for access/monitoring). Oesophageal Doppler typically unhelpful.
  \item OLV facilitates surgery and prevents soiling.
  \item Airway pressure/volume loop display is a useful adjunct.
  \item Surgical manipulation can reduce CO/BP and cause arrhythmias.
  \item Suction collapsed lung prior to reinflation.
  \item Bronchial suture line leak test under saline with manual inflation to high pressure (as described).
  \item Fluids: titrate to losses; avoid excessive replacement, particularly pneumonectomy.
  \item Regional catheter analgesia (epidural/paravertebral); cautious incremental dosing if established pre-op.
\end{itemize}

\subsection{7.3 Postoperative}
\begin{itemize}[leftmargin=*, itemsep=2pt]
  \item Aim to extubate awake and sitting.
  \item Supplementary oxygen (humidified preferable; nasal cannulae more likely to remain in place on ward).
  \item Ensure good analgesia.
  \item CXR usually required in recovery.
\end{itemize}

\subsection{7.4 Special considerations}
\begin{itemize}[leftmargin=*, itemsep=2pt]
  \item Non-metastatic manifestations of bronchial carcinoma may occur (e.g. paraneoplastic/neuromuscular syndromes).
  \item Pneumonectomy: notable perioperative risks including acute lung injury; consider risk factors and protective strategies.
  \item Arrhythmias (especially AF) are common post-pneumonectomy; prophylaxis strategies may be used according to local practice.
\end{itemize}

\section{8. Thoracoscopy / VATS procedures}

\subsection{8.1 Preoperative}
Similar assessment to thoracotomy; generally less invasive with less postoperative deterioration in lung function. Discuss analgesia plan (regional/PCA) where appropriate.

\subsection{8.2 Perioperative}
\begin{itemize}[leftmargin=*, itemsep=2pt]
  \item Consider arterial line in high-risk/compromised patients.
  \item Upper arm IV; radial arterial line usually in dependent arm.
  \item Intra-op opioid boluses.
  \item Commence OLV (left DLT) before trocar insertion; need good lung collapse.
  \item Intercostal or paravertebral blocks; paravertebral catheter can be placed thoracoscopically for more extensive procedures.
\end{itemize}

\subsection{8.3 Postoperative}
\begin{itemize}[leftmargin=*, itemsep=2pt]
  \item Extubate, sit up, start supplementary oxygen before transfer.
  \item CXR in recovery to confirm full re-expansion.
  \item Balanced analgesia; PCA morphine may be required for more painful procedures.
  \item Encourage early mobilisation.
\end{itemize}

\subsection{8.4 Special considerations}
Always possible conversion to open thoracotomy. Epidural not usually necessary; consider if bilateral.

\section{9. Lung volume reduction surgery (LVRS) and bullectomy}

\subsection{9.1 Concept}
Surgical treatment for selected severe emphysema; aims to reduce total lung volume by resecting most diseased areas \(\rightarrow\) improve function. Patient group often very high risk for GA.

\subsection{9.2 Preoperative}
Intensive assessment, careful selection, optimisation. Cardiac assessment may include coronary angiography + right heart catheterisation. Steroid supplementation often required. Requires adequate thoracic experience and understanding of pathophysiology.

\subsection{9.3 Perioperative}
Approach may be sternotomy/thoracotomy/VATS. Consider protective ventilation strategies.

\subsection{9.4 Postoperative}
HDU/ICU often required; extubate ASAP. Accept higher PaCO\(_2\) within target range and adjust \(\mathrm{FiO_2}\) to maintain acceptable saturation. Watch for air leaks; limit drain suction. Requires excellent analgesia, skilled physiotherapy, pulmonary rehabilitation.

\subsection{9.5 Special considerations}
Commonest complication: prolonged air leak. Notable mortality in reported series. LVRS benefits particular emphysema phenotypes (as described in source). Isolated congenital bulla/cyst patients may be fitter but still need careful intra-op management.

\section{10. Drainage of empyema and decortication}

\subsection{10.1 Preoperative}
Intrapleural infection often secondary to pneumonia, drains, or chest surgery. Patients may be debilitated/septic. Respiratory function may be compromised by pneumonia or previous resection. Check for bronchopleural fistula.

\subsection{10.2 Perioperative}
\begin{itemize}[leftmargin=*, itemsep=2pt]
  \item Empyema may be drained by rib resection and large-bore drain; thoracoscopy may break loculations/adhesions.
  \item Decortication requires ``thoracotomy'' anaesthetic with epidural analgesia (paravertebral often not possible due to loss of pleura).
  \item Decortication can cause significant haemorrhage.
  \item Consider arterial line/CVP monitoring for all but the fittest.
\end{itemize}

\subsection{10.3 Postoperative}
Balanced analgesia; HDU care recommended for debilitated patients.

\subsection{10.4 Special considerations}
Surgical principle: remove infected tissue/pleural ``peel'', re-expand lung, obliterate space. Air leaks common after visceral pleura decortication; lobectomy may be required for massive leaks/severe parenchymal damage.

\section{11. Repair of bronchopleural fistula (BPF)}

\subsection{11.1 Preoperative}
Features: productive cough, haemoptysis, fever, dyspnoea, subcutaneous emphysema, falling fluid level in post-pneumonectomy space on CXR. Severity proportional to fistula size; large fistulae may need urgent support. Often debilitated with infection and reduced reserve. Review prior anaesthetic charts and airway anatomy. Ensure oxygen, functioning chest drain, IV antibiotics, fluids.

\subsection{11.2 Perioperative: core principles}
Protect the good lung from contamination and control ventilation distribution. Failure to isolate lungs after induction carries grave risk. Small/moderate fistulae: bronchoscopy \(\pm\) tissue glue. Start invasive arterial monitoring before induction.
\begin{itemize}[leftmargin=*, itemsep=2pt]
  \item Airway strategy described: awake intubation under LA often described as safest.
  \item Many thoracic anaesthetists use modified RSI and advance DLT under fibreoptic guidance to bronchus contralateral to fistula.
  \item Risk: fistula may be enlarged by inappropriate DLT placement.
\end{itemize}

\subsection{11.3 Postoperative}
Plan HDU/ICU except straightforward cases. Minimise airway pressures; extubate ASAP. Use post-thoracotomy analgesia and monitor renal function if NSAIDs used.

\subsection{11.4 Special considerations}
Usually post-pneumonectomy/lobectomy; can be secondary to pneumonia, lung abscess, empyema. Anaesthesia challenging; not recommended for occasional thoracic practice.

\subsection{11.5 Tips for controlling a massive air leak (unable to ventilate effectively)}
If DLT cannot be positioned satisfactorily:
\begin{itemize}[leftmargin=*, itemsep=2pt]
  \item Intubate with uncut cuffed 6 mm single-lumen tube; fibreoptic into intact main bronchus and railroad tube into bronchus.
  \item Surgeon passes rigid bronchoscope into intact main bronchus; slide long flexible bougie/airway exchange catheter (jet ventilation possible) then railroad single-lumen tube.
  \item If all else fails: bronchial blocker (e.g. Arndt) or large Fogarty catheter into fistula via rigid bronchoscope to control leak temporarily.
\end{itemize}

\section{12. Pleurectomy / pleurodesis}

\subsection{12.1 Preoperative}
Two typical groups: (1) young/fit with recurrent pneumothoraces (check asthma), (2) older with COPD or recurrent pleural effusions (check reserve). Large unilateral effusions rarely cause orthopnoea if otherwise healthy; orthopnoea suggests additional pathology (e.g. heart failure). Check recent CXR. Pre-op intercostal drain advised if pneumothorax present. Confirm approach and discuss postoperative analgesia/regional technique.

\subsection{12.2 Perioperative}
\begin{itemize}[leftmargin=*, itemsep=2pt]
  \item Keep airway pressures low if pneumothorax history; avoid nitrous oxide.
  \item Be alert for pneumothorax/tension pneumothorax during IPPV (can occur despite drains and on ``healthy'' side).
  \item Instillation of irritant: collapse lung (if single-lumen tube, preoxygenate then briefly disconnect).
  \item End: aim for full lung expansion to appose parietal and visceral pleura.
\end{itemize}

\subsection{12.3 Postoperative}
Extubate and sit upright; CXR to confirm full expansion; drain suction often prescribed.

\subsection{12.4 Pain and analgesia notes}
Pleural inflammation is painful, especially with abrasion. Regular paracetamol. Avoid NSAIDs if they may reduce pleurodesis effectiveness. Thoracic epidural recommended for pleurectomy (especially bilateral); PCA morphine + intercostal blocks as alternative. Paravertebral blocks often unsuitable due to pleural damage.

\subsection{12.5 Special considerations}
Pleurectomy often for recurrent pneumothorax with stapling of leaking tissue (apical blebs/bullae). Pleurodesis commonly for malignant effusions. Massive effusions: consider partial drainage \(\geq 12\) h pre-op to reduce risk of re-expansion pulmonary oedema from rapid reinflation. Risk of circulatory collapse when turned ``effusion side up'' (mediastinal shift + high intrathoracic pressure reducing venous return/CO): return supine and drain before proceeding.

\section{13. Oesophagectomy}

\subsection{13.1 Preoperative}
Usually oesophageal cancer; occasionally benign stricture/achalasia. Understand approach (transhiatal; Ivor--Lewis; thoracoabdominal; McKeown 3-stage; minimally invasive combinations). Malnutrition/cachexia common \(\rightarrow\) increased morbidity/mortality; careful cardiorespiratory assessment. Plan for long duration + potential repositioning. Aspiration risk: ranitidine/omeprazole if swallowable. Arrange HDU/ICU depending on protocols/fitness.

\subsection{13.2 Perioperative}
\begin{itemize}[leftmargin=*, itemsep=2pt]
  \item Treat as aspiration risk: RSI with cricoid pressure advised.
  \item If thoracotomy planned: DLT + OLV to facilitate access and reduce lung trauma.
  \item Regional technique by approach:
  \begin{itemize}[leftmargin=*, itemsep=2pt]
    \item Thoracoabdominal: paravertebral LA infusion + morphine PCA.
    \item Laparotomy/thoracotomy: mid-thoracic epidural (incremental dosing and post-op infusion).
  \end{itemize}
  \item NG tube required (removed for resection, reinserted under surgical guidance after anastomosis).
  \item Avoid IJV line on side needed for cervical anastomosis.
  \item Temperature control: core temperature monitoring and active warming (fluid warmer + forced-air warming).
  \item Fluids: stay ahead; monitor Hb and blood gases; watch metabolic acidosis.
  \item Intrathoracic mobilisation may provoke arrhythmias and reduced CO with hypotension.
  \item Change DLT to single-lumen tube before cervical anastomosis (if performed) to improve access.
\end{itemize}

\subsection{13.3 Postoperative}
Requires intensive experienced nursing in specialist ward/HDU/ICU. If hypothermic (\(\textless 35.5^\circ\)C) or unstable: ventilate until improved. Aim minimum urine output 1 mL/kg/h. Early enteral feeding via jejunostomy or nasoduodenal tube.

\subsection{13.4 Special considerations}
High elective perioperative mortality (even in specialist centres). Major complication rate high; deaths often due to systemic sepsis from respiratory complications or anastomotic breakdown. Minimally invasive approaches increasing; beware tension capnothorax if pleura breached during laparoscopic hiatal dissection. Not recommended as occasional practice.

\section{14. Chest injury (thoracic trauma)}

\subsection{14.1 Context}
Initial diagnosis and management of major thoracic trauma described elsewhere. Here: anaesthetic management for definitive repair (diaphragm, oesophagus, tracheobronchial tree).

\subsection{14.2 General considerations}
\begin{itemize}[leftmargin=*, itemsep=2pt]
  \item Often with major head/abdominal/skeletal injuries \(\rightarrow\) prioritise appropriately (C-spine, laparotomy for bleeding, splintage).
  \item \(\textless 30\%\) require thoracotomy, but persistent bleeding from drains \(>200\) mL/h indicates urgent surgery.
  \item Most deaths due to exsanguination \(\rightarrow\) ensure rapid infusion capability (two large-bore IVs).
  \item Maintain high suspicion for tension pneumothorax during IPPV (drain does not guarantee).
  \item Massive air leaks suggest significant tracheobronchial injury.
  \item High risk of multiple organ failure \(\rightarrow\) postoperative ICU care.
\end{itemize}

\subsection{14.3 Repair of ruptured diaphragm}
May present chronically or as bowel obstruction \(\rightarrow\) correct fluid/electrolytes. Repair should be prompt but seldom true emergency. Approach: lateral thoracotomy or thoracoabdominal incision. Management as for fundoplication. Avoid nitrous oxide (bowel distension). DLT + OLV facilitate access. NG tube to decompress stomach.

\subsection{14.4 Repair of ruptured oesophagus}
Diagnosis described elsewhere; surgical emphysema and pleural effusions often present. Other causes include Boerhaave's syndrome, foreign bodies, iatrogenic injury.

\subsection{14.5 Tracheobronchial injury: definitive repair (key anaesthetic points)}
\begin{itemize}[leftmargin=*, itemsep=2pt]
  \item Priority: \(100\%\) \(\mathrm{O_2}\) and relieve tension pneumothorax (may need two large-bore drains with independent underwater seals).
  \item If ventilation/oxygenation acceptable: call thoracic help and assess site with fibreoptic bronchoscopy before intubation.
  \item Principles as for large BPF: adequate PPV may be impossible with single-lumen tube.
  \item Isolate torn bronchus by fibreoptic-guided intubation of contralateral intact main bronchus with appropriate DLT.
  \item Alternatively, guide an uncut single-lumen tube past an upper tracheal tear so cuff lies distal to injury.
  \item Once airway secure: urgent thoracotomy for repair.
  \item Carinal disruption may require cardiopulmonary bypass for oxygenation during repair.
  \item Poor management can lead to later stenosis and long-term airway problems.
\end{itemize}

\section{15. Procedure summaries (quick reference)}

\subsection{15.1 Fibreoptic bronchoscopy}
Visual inspection of tracheobronchial tree \(\pm\) biopsy/brushings/lavage. Time 5--10 min; pain \(+\); supine; blood loss none. GA rarely used. Single-lumen tube (8--9 mm) with bronchoscopy diaphragm on angle piece; IPPV with relaxant appropriate to duration. Expect high airway pressures while scope in ETT; suction can empty breathing system.

\subsection{15.2 Lung biopsy}
Diagnostic sampling (localised or diffuse). Time 30--60 min; pain \(+++/++++\). Lateral, VATS or minithoracotomy. Blood loss minor; group \& save (crossmatch if anaemic). DLT + OLV facilitate VATS. Diffuse disease patients may have very poor lung function with risk ventilator dependence and significant mortality.

\subsection{15.3 Oesophagoscopy and dilatation}
Inspection via rigid/fibreoptic scope \(\pm\) dilatation (bougies or balloon). Core issues: regurgitation/aspiration risk and shared airway.

\subsection{15.4 Rigid bronchoscopy}
Inspection \(\pm\) biopsy, stents, foreign body removal. Time 5--20 min; pain \(+\); supine with head/neck extended; blood loss usually minimal. Technique note: TIVA; short-acting opioid; intermittent suxamethonium; ventilation via bronchoscope (Venturi/Sanders injector).

\subsection{15.5 Superior/cervical mediastinoscopy}
Inspect/biopsy superior/anterior mediastinum via suprasternal/anterior intercostal incision. Time 20--30 min; pain \(+\); supine/slightly head-up. Blood loss usually minimal but potential massive haemorrhage.

\subsection{15.6 Lung resection (wedge/lobectomy/pneumonectomy)}
Time 2--4 h; pain \(+++++\). Lateral, table ``broken''. Blood loss 200--800 mL (can be more). Technique: IPPV via DLT with OLV; regional catheter for analgesia; arterial line for pneumonectomy/less fit.

\subsection{15.7 VATS procedures}
Time 30--120 min; pain \(++/+++\). Lateral, table ``broken''. Blood loss minimal--200 mL. Technique: IPPV + OLV via left DLT; intercostal/paravertebral blocks \(\pm\) catheter; consider arterial line.

\subsection{15.8 Lung volume reduction surgery / bullectomy}
Core issues: severe respiratory compromise, air leak risk, postoperative ventilatory strategy.

\subsection{15.9 Empyema drainage / decortication}
Drainage time 20--40 min; decortication 2--3 h. Pain \(+++/+++++\). Blood loss: drainage minimal; decortication 500--2000 mL (crossmatch often required). Technique: GA; DLT advised for decortication (air leak risk); arterial line/CVP.

\subsection{15.10 Bronchopleural fistula repair}
Time 2--3 h (thoracotomy). Pain \(++++/+++++\). Keep sitting upright with affected side down until good lung isolated, then lateral. Blood loss 300--800 mL. Technique: fibreoptic-guided endobronchial intubation with DLT (various induction strategies described).

\subsection{15.11 Pleurectomy / pleurodesis}
Pleurectomy 1--2 h; pleurodesis 20--40 min. Pain \(+++/++++\). Blood loss minimal thoracoscopic; up to \(\sim 500\) mL for thoracotomy. Technique: IPPV, DLT + OLV for open/VATS; single-lumen tube often adequate for talc pleurodesis.

\subsection{15.12 Oesophagectomy}
Time 3--6 h; pain \(+++++\). Position: supine \(\pm\) lateral for thoracotomy. Blood loss 500--1500 mL (crossmatch often required). Technique: IPPV; DLT useful if thoracotomy; arterial/CVP lines; urinary catheter; thoracic epidural or paravertebral catheter.

\section{16. Addendum: Blindspots in Oxford Ch15 (filled from other texts)}

\subsection{A1. HPV and why oxygenation can improve after several minutes of OLV}
Oxford describes shunt and practical manoeuvres, but is relatively light on the physiology of hypoxic pulmonary vasoconstriction (HPV) and how anaesthesia interacts with it.

\begin{itemize}[leftmargin=*, itemsep=2pt]
  \item \textbf{Volatile agents inhibit HPV mainly at higher doses:} for volatile agents, the ED\(_{50}\) for inhibition of HPV is about \textbf{2 MAC} (i.e., substantial inhibition occurs at higher concentrations).\\
  \emph{Source: Morgan \& Mikhail (5e), Ch 23 Respiratory Physiology \& Anesthesia; PDF p. 524.}
\end{itemize}

\subsection{A2. Device selection: why some clinicians avoid bronchial blockers (exam-viva fodder)}
Oxford is pragmatic about DLT vs blockers; Morgan \& Mikhail explicitly spells out the common disadvantages that become MCQ/viva “gotchas”.

\begin{itemize}[leftmargin=*, itemsep=2pt]
  \item \textbf{Compared with DLTs, bronchial blockers:}
  \begin{itemize}[leftmargin=*, itemsep=2pt]
    \item have \textbf{increased dislodgement},
    \item have \textbf{small central lumens} that \textbf{do not allow efficient suctioning} of secretions,
    \item and may not allow \textbf{rapid lung collapse} (collapse may be slow/incomplete because of the small channel).
  \end{itemize}
  \emph{Source: Morgan \& Mikhail (5e), Ch 25 Anesthesia for Thoracic Surgery; PDF p. 568.}
\end{itemize}

\subsection{A3. DLT-specific complications that Oxford under-emphasises}
Oxford focuses on placement/confirmation; Morgan \& Mikhail provides a compact list of DLT-specific complications that are commonly examined.

\begin{itemize}[leftmargin=*, itemsep=2pt]
  \item \textbf{Complications include:}
  \begin{itemize}[leftmargin=*, itemsep=2pt]
    \item \textbf{laryngitis},
    \item \textbf{tracheobronchial rupture} from traumatic placement or \textbf{overinflation of the bronchial cuff},
    \item \textbf{inadvertent suturing} of the tube to a bronchus (suggested by inability to withdraw the tube at extubation).
  \end{itemize}
  \emph{Source: Morgan \& Mikhail (5e), Ch 25 Anesthesia for Thoracic Surgery; PDF p. 568.}
\end{itemize}

\subsection{A4. Postoperative: structured care bundle + haemorrhage trigger}
Oxford covers CXR/O\(_2\)/analgesia but is less “bundle + thresholds”. Morgan \& Mikhail gives a very exam-friendly set of routine items and a numerical haemorrhage trigger.

\begin{itemize}[leftmargin=*, itemsep=2pt]
  \item \textbf{Postoperative haemorrhage warning signs:} increased chest tube drainage \textbf{>200 mL/h} with hypotension, tachycardia, and falling haematocrit.
  \item \textbf{Routine postoperative care bundle:}
  \begin{itemize}[leftmargin=*, itemsep=2pt]
    \item semiupright position \textbf{>30\(^\circ\)},
    \item supplemental oxygen \textbf{40--50\%},
    \item \textbf{incentive spirometry},
    \item ECG and haemodynamic monitoring,
    \item postoperative chest radiograph (confirm drain/line position and lung expansion),
    \item adequate pain relief.
  \end{itemize}
  \emph{Source: Morgan \& Mikhail (5e), Ch 25 Anesthesia for Thoracic Surgery; PDF p. 575.}
\end{itemize}

\subsection{A5. Bullae/cysts: induction strategy + nitrous oxide (classic trap)}
Oxford mentions avoidance of nitrous oxide in relevant contexts; Morgan \& Mikhail makes the “why + how” explicit and links to tension physiology.

\begin{itemize}[leftmargin=*, itemsep=2pt]
  \item In patients with cysts/bullae, \textbf{induction with maintenance of spontaneous ventilation} is described as desirable until the affected side is isolated with a DLT (or a chest tube is placed), because rupture during PPV can cause \textbf{tension pneumothorax}.
  \item \textbf{N\(_2\)O is contraindicated} because it can expand the air space and precipitate rupture.
  \emph{Source: Morgan \& Mikhail (5e), Ch 25 Anesthesia for Thoracic Surgery; PDF p. 577.}
\end{itemize}

\end{document}
